\Needspace{8\baselineskip}\originalchapter{习题课：核素、探测与加速器}
本章把 Recitation 16--20 的十道原题放在同一条测量链上：质谱给质量，质量差给核能，衰变与径迹给时间及能损，计时与同步辐射限制加速器。题面数字按原件保留。官方扫描解答的思路与本稿重新计算或修正之处逐题注明。
\originalsection{质量测定和核能}
\begin{exercise}\label{ex:rec-16-1}
Recitation 16题1。核结合能须从准确原子质量求得。说明20世纪初发展起来的质谱法如何利用电场、磁场中带电量$Q$离子的偏转，同时测得其动量、动能和质量。
\end{exercise}
\begin{solution}
官方方法重编。先使离子从静止跨过已知电势差$V$，则$K=|Q|V$；若不是从静止出发，须测入射能并加上$|Q|V$。在均匀横磁场中量轨迹半径$r_B$，洛伦兹力$|Q|vB=pv/r_B$给$p=|Q|Br_B$。非相对论区域
\begin{equation}M=\frac{p^2}{2K}=\frac{|Q|B^2r_B^2}{2V}.
\end{equation}
另一种独立读出$K$的方法是让离子在已知径向电场$E$中转过半径$r_E$，由$Mv^2/r_E=|Q|E$得$2K=|Q|Er_E$。两场区的动能若被电极改变，需记录各处电势再对应，不能无条件视作同一个速度。原扫描解把磁半径一处写成$r_E$，这里分别标$r_B,r_E$。在相对论区域仍有$p=|Q|Br_B$，但应使用$M^2c^4=E_{\rm tot}^2-p^2c^2$而非$K=p^2/(2M)$。获得中性原子质量还要依电离态校正电子质量与束缚能；再以质量亏损求核结合能。
\end{solution}
\begin{exercise}\label{ex:rec-16-2}
Recitation 16题2。$^{238}$Pu的$\alpha$衰变平均寿命$\tau=127$年，生成长寿命$^{234}$U（平均寿命$3.5\times10^5$年），每次释放$5.49$MeV动能。旅行者2号在1977年8月20日发射、1981年8月26日到土星、1989年8月24日到海王星。若热电转换效率为$5.5\%$，为到达土星时至少输出395W，出发时需带多少钚？到达海王星时还能输出多少电功率？
\end{exercise}
\begin{solution}
官方指数衰变方法重编，日期和题面能量按一致精度复算。题给$\tau$是平均寿命，因此$N(t)=N_0\e^{-t/\tau}$，不是把127年代入半衰期式。以题设把每次5.49MeV作为可收集热量、转换效率恒定，则
\begin{equation}P_e(t)=\eta\frac{N_0}{\tau_s}Q_\alpha\e^{-t/\tau},\qquad
N_0=\frac{P_e(t_S)\tau_s\e^{t_S/\tau}}{\eta Q_\alpha}.
\end{equation}
$Q_\alpha=5.49\times10^6(1.602176634\times10^{-19})$J、$\tau_s=127(365.25)(86400)$s，$t_S=4.01643$年。代入得$N_0=3.37749\times10^{25}$、初始钚质量$N_0(238\,\mathrm{g\,mol^{-1}})/N_A\simeq13.35$kg。海王星比土星晚$7.99452$年，故$P_N=395\e^{-7.99452/127}\simeq370.9$W。官方粗算约13.3kg、370W。铀子体在这一时间尺度的衰变热忽略。若另计$\alpha$反冲热、热电器老化，数值会改变；那是题设以外的工程模型。
\end{solution}
\originalsection{结合能、放射性测年与壳层}
\begin{exercise}\label{ex:rec-17-1}
Recitation 17题1。$^{56}_{26}$Fe的原子质量为$55.9349$u。求每核子结合能（MeV），并说明原子质量代入时电子质量如何处理。
\end{exercise}
\begin{solution}
官方质量亏损法重编，改用一致精度的质量。用氢原子质量$m_H$而非裸质子质量，原子两侧的26个电子相消，忽略小的电子束缚能修正：
\begin{equation}
\frac BA=\frac{[26m_H+30m_n-m(^{56}\mathrm{Fe})]c^2}{56}.
\end{equation}
取$m_H=1.00782503223$u、$m_n=1.00866491595$u、$1\mathrm u c^2=931.49410242$MeV，得$B\simeq492.293$MeV，$B/A\simeq8.79095$MeV。官方扫描以粗舍入质子、中子和电子质量得到约8.82MeV；两种数字精度不同，不能把后者与上述精确输入混算。
\end{solution}
\begin{exercise}\label{ex:rec-17-2}
Recitation 17题2。陶片含碳1g，现测$^{14}$C活度$0.0231$Bq。生物活着时$^{14}$C/$^{12}$C约$1.2\times10^{-12}$，$^{14}$C半衰期5730年。估计陶片年代。
\end{exercise}
\begin{solution}
官方放射性测年法重编，并改正解答重印题面时把$10^{-12}$错写成$10^{12}$。近似把1g碳视为$^{12}$C，初始$^{14}$C数
$N_0=(1/12)N_A(1.2\times10^{-12})\simeq6.02214\times10^{10}$。
$\lambda=\ln2/[5730(365.25)(86400)]\simeq3.83325\times10^{-12}\mathrm{s}^{-1}$，故初活度$A_0=\lambda N_0\simeq0.230844$Bq。
\begin{equation}t=\frac{5730}{\ln2}\ln\frac{A_0}{0.0231}\simeq1.903\times10^4\ \mathrm{yr}.
\end{equation}
模型前提是陶器原料停止与环境交换碳，初始同位素比适用且污染/本底已扣除。若材料是经复杂烧制与后期附着的陶片，只测总碳不自动满足这些条件。
\end{solution}
\begin{exercise}\label{ex:rec-18-1}
Recitation 18题1。$^{250}_{96}$Cm的质量题给为$232.938$GeV，原件未注明核质量或原子质量口径。根据$A,Z$判断是否预期稳定，并说明这些信息能和不能判定哪些衰变。
\end{exercise}
\begin{solution}
官方解用$Z/A=96/250=0.384$和重核经验值约0.41比较，预期这种极重核不会永久稳定；但单凭固定比例不能判定具体衰变道。题给质量若不先与核/原子质量口径一致的其他质量比较，也不能单独给未知子体的$Q$值。实际上同卷液滴模型习题\ref{ex:8701-5-2}在$A=250$处给连续$\beta$谷$Z_*\simeq96.65$，Cm的$Z=96$就在谷附近。因此“比例偏离0.41所以必$\beta$不稳定”并不成立。极重核还须逐一比较$\alpha$衰变、自发裂变和可及$\beta$末态的质量及势垒；题面没有所有子体质量与半衰期，不能由这一个数推出主导机制。这里把“预期不永久稳定”与“已证实某道放能”分开。
\end{solution}
\begin{exercise}\label{ex:rec-18-2}
Recitation 18题2。若核力电荷无关且一个中子、一个质子能形成束缚氘核，为何两个中子或两个质子没有相应的束缚二核子态？这告诉我们核子间作用的什么性质？
\end{exercise}
\begin{solution}
官方Pauli思想重编。低能$S$波的空间波函数对交换对称，总费米子态须反对称。氘核是$S=1,I=0$的自旋三重态、同位旋单重态（还含小的${}^3D_1$混合）；$nn$或$pp$同种核子的同位旋只能$I=1$，其低能$S$波必须是自旋单重${}^1S_0$。电荷无关仅说在同一$(S,I,L)$道旋转质子/中子时强相互作用近似相同，并不强制${}^1S_0$与${}^3S_1$有同样束缚能。${}^1S_0$虽强吸引却不束缚；$pp$还受库仑排斥。结论是核力依赖自旋与同位旋通道，不能笼统说“任意两个相同核子都不能在高角动量束缚”。
\end{solution}
\begin{exercise}\label{ex:rec-18-3}
Recitation 18题3。简述促成核壳层模型的主要实验现象；给出闭壳核实例，并明确各例是哪一种核子数闭壳。
\end{exercise}
\begin{solution}
官方魔数列表重编并补足原解遗漏的实例。核子分离能、第一激发能和稳定性随$N$或$Z$变化在$2,8,20,28$，以及$50,82,126$附近出现不连续增强，提示平均场能级充满后新增核子进入新壳；自旋轨道耦合解释28以上魔数。例子：$^4$He的$Z=N=2$、$^{16}$O的$Z=N=8$、$^{40}$Ca的$Z=N=20$均双闭壳；$^{48}$Ca的$Z=20,N=28$，$^{208}$Pb的$Z=82,N=126$也双闭壳。魔数是现象归纳，不是把所有这些核都宣称绝对不衰变。
\end{solution}
\originalsection{云室、计时与辐射损失}
\begin{exercise}\label{ex:rec-19-1}
Recitation 19题1。Anderson的1931年云室正电子径迹（原题标题误称泡室）所在磁场$B=1.5$T、指向纸内；粒子从下方入射，穿过6mm厚铅板。板前后曲率半径分别15.5cm、5.3cm。(a)求前后动量及电荷符号；(b)用题给近似Bethe式$-dE/dX=4\pi N_A r_e^2m_ec^2z^2(Z/A)\beta^{-2}\ln(m_ec^2\gamma^2\beta^2/I)$，将铅取$Z=82,A=207,I=820$eV、$\rho=11.34\mathrm{g\,cm^{-3}}$，以恒定入射阻止本领估计质子、带电pion、电子的6mm电离损失，判断穿越者身份；(c)再按题给铅的辐射长度$X_0=0.56$cm，积分$\dd E/\dd x=-E/X_0$比较轫致辐射。给定$N_A=6.022\times10^{23}\mathrm{mol}^{-1}$，$m_ec^2=511$keV，$r_e=e^2/(4\pi\epsilon_0m_ec^2)$。
\end{exercise}
\begin{center}\begin{tikzpicture}[>=Stealth,scale=.85]
\draw[fill=gray!20] (-2,-.06) rectangle (2,.06);
\node[right] at (2,.05) {铅板，6mm};
\draw[->,very thick] (.35,-1.6) .. controls (.35,-1) and (.06,-.45) .. (-.4,-.07);
\draw[->,very thick] (-.4,.07) .. controls (-.78,.4) and (-1.3,.6) .. (-1.9,.68);
\node[right] at (.35,-1.1) {自下入射，$r_i=15.5$cm};
\node[above] at (-1.55,.68) {$r_f=5.3$cm};
\node at (1.1,1.05) {$\otimes\,\mathbf B$指入纸面};
\end{tikzpicture}\\[-2pt]
{\small 根据原公共领域照片的弯曲方向自绘示意；曲率不按比例。}
\end{center}
\begin{solution}
官方辨认思路重编，能损近似的失效处由编者标明。对$|Q|=e$，$p[\mathrm{GeV}/c]=0.299792458B[\mathrm T]r[\mathrm m]$，得$p_i\simeq69.70$、$p_f\simeq23.83$MeV$/c$；按题给入射方向、弯曲方向和指入纸面的磁场判为正电荷（参见原件公共领域照片链接）。仅有两个曲率半径数值而无弯曲朝向，数学上本不能判符号；这里符号使用了原图的方向信息。源解把半径的m误写cm，且把铅密度单位误写为g/cm$^2$，此处已更正。

质量厚度$X=\rho x=11.34(0.6)=6.804\mathrm{g\,cm^{-2}}$。按入射$p_i$分别算$\gamma\beta=p_i/(mc)$，把$-dE/dX$乘$X$，得到固定本领的电离估计约质子186MeV、pion21MeV、电子13.5MeV。质子初动能仅2.59MeV，pion约16.4MeV；两者的恒本领估计都大于可用动能，说明必须沿轨迹更新速度且会停止，不能把负的“末能”当物理结果。电子初能约69.7MeV，较轻粒子可穿过；不过题给Bethe式本身并非精确的电子阻止本领公式。

轫致辐射单独求解得$E(x)=E_i\e^{-x/X_0}$，$x=0.6$cm时$E_f\simeq23.9$MeV，和曲率所示末动量约23.8MeV同量级。若把固定电子电离本领$S\simeq22.44$MeV/cm和辐射同时放进平均能量方程$E'=-E/X_0-S$，则$E_f=(E_i+SX_0)\e^{-x/X_0}-SX_0\simeq15.6$MeV，而非官方把先后两个平均损失直接相减得到的约11MeV。因此单条随机径迹不能据这些平均模型精确等同；能损涨落和实际能量依赖均需更完整处理。$X$以g/cm$^2$计，辐射式的$x,X_0$以cm计，不可混用。
\end{solution}
\begin{exercise}\label{ex:rec-20-1}
Recitation 20题1。两个质量$m_1,m_2$、动量同为$p$的粒子跨越距离$L$，求飞行时间差。两个闪烁计数器对时间差的分辨率$\sigma_{\Delta t}=300$ps；要以$2\sigma$分辨$p=4$GeV$/c$的$\pi$与$K$，间距至少多大？
\end{exercise}
\begin{solution}
官方计算重编，时间差取正的$t_K-t_\pi$：
\begin{equation}
\Delta t=\frac Lc\left[\sqrt{1+\frac{m_K^2c^2}{p^2}}-\sqrt{1+\frac{m_\pi^2c^2}{p^2}}\right].
\end{equation}
令$\Delta t=2(300\mathrm{ps})$，代$m_\pi c^2=140$MeV、$m_Kc^2=494$MeV，得到$L\simeq25.75$m。官方扫描在一个中间差式反写次序造成负号，最终距离仍取正。若300ps指每台计数器各自独立的时间分辨率，差的标准差会是$\sqrt2$倍、距离约36.4m；但那不是本题指定的差测量分辨率。
\end{solution}
\begin{exercise}\label{ex:rec-20-2}
Recitation 20题2。圆形对撞机曲率半径都取$R=4.3$km；电子--正电子质心总能量200GeV、质子--质子质心总能量14TeV。求单粒子每圈同步辐射能量损失，分别给电子和质子结果。
\end{exercise}
\begin{solution}
官方Liénard公式方法重编。每束能量分别是$E_e=100$GeV、$E_p=7$TeV，不是把质心总能量代入单粒子$\gamma$。对$\beta\simeq1$的圆周运动，辐射功率$P=q^2\gamma^4a^2/(6\pi\epsilon_0c^3)$与周长/速度相乘，得到
\begin{equation}U_{\rm turn}=\frac{q^2\gamma^4}{3\epsilon_0R}=\frac{4\pi\alpha\hbar c}{3R}\gamma^4.
\end{equation}
以$m_ec^2=0.510999$MeV、$m_pc^2=938.272$MeV，得电子约2.057GeV/圈、质子约4.35keV/圈。差异主要来自固定能量下$\gamma=E/(mc^2)$的四次方。题设把整圈视为相同曲率的理想环，未逐段建模真实直线区与磁铁弯曲半径。
\end{solution}
